\item \textbf{Problema de repaso.}
Una roca descansa sobre una acera de concreto. Se presenta un terremoto, que mueve al suelo verticalmente en movimiento armónico con una frecuencia constante de $2.40\text{ Hz}$ y con amplitud gradualmente creciente.
\begin{enumerate}
    \item ¿Con qué amplitud vibra el suelo cuando la roca comienza a perder contacto con la acera?
\end{enumerate}
Otra roca está asentada sobre el concreto en el fondo de una alberca llena con agua. El terremoto sólo produce movimiento vertical, así que el agua no salpica de lado a lado.
\begin{enumerate}
    \setcounter{enumi}{1}
    \item Presente un argumento convincente de que, cuando el suelo vibra con la amplitud encontrada en el inciso (1), la roca sumergida también apenas pierde contacto con el suelo de la alberca.
\end{enumerate}